\chapter{de Rham 预叠}
\markboth{第 1 章\ de Rham 预叠}{第 1 章\ de Rham 预叠}

对预叠 $Y$，晶体定义为 $Y$ 的 de Rham 预叠 $Y_{\dR}$ 上的层（拟凝聚或归纳凝聚）。本节定义函子 $Y\mapsto Y_{\dR}$，并建立它的一些基本性质。

最重要的是，我们将证明：若 $Y$ 局部几乎有限型，则 $Y_{\dR}$ 也局部几乎有限型。此时我们还将证明 $Y_{\dR}$ 是经典的，即可以完全在「经典」代数几何的范围内研究它，无须涉及导出环。

读者可能觉得本节格外抽象。一个可行的策略是初读时跳过本节，等到本节建立的论断用于晶体时再回头查阅。

\section{定义与基本性质}

\ob{1.1.1}{n:1.1.1}设 $Y$ 是 $\PreStk$ 的对象。定义 $Y$ 的 de Rham 预叠 $Y_{\dR}\in\PreStk$ 为
\begin{equation}
Y_{\dR}(S):=Y(\clred S),\qquad S\in\DGSch_{\mathrm{aff}}.\tag{1.1}
\end{equation}

\ob{1.1.2}{n:1.1.2}更抽象地，可以改写为
\[
Y_{\dR}:=\mathrm{RKE}_{{}^{\mathrm{red}}\Sch_{\mathrm{aff}}\hookrightarrow\DGSch_{\mathrm{aff}}}(\clred Y),
\]
其中 $\clred Y:=Y|_{{}^{\mathrm{red}}\Sch_{\mathrm{aff}}}$ 是 $Y$ 到既约经典仿射概形的限制，而
\[
\mathrm{RKE}_{{}^{\mathrm{red}}\Sch_{\mathrm{aff}}\hookrightarrow\DGSch_{\mathrm{aff}}}
\]
是函子沿包含 ${}^{\mathrm{red}}\Sch_{\mathrm{aff}}\hookrightarrow\DGSch_{\mathrm{aff}}$ 的右 Kan 延拓。

\ob{1.1.3}{n:1.1.3}下面这个（显然的）观察在后续会用到。

\begin{othm}{引理}{1.1.4}{n:1.1.4}
函子 $\dR:\PreStk\to\PreStk$ 与极限、余极限交换。
\end{othm}

\begin{proof}
由 $\PreStk=\Funct\big((\DGSch_{\mathrm{aff}})^{\mathrm{op}},\infty\text{-}\Grpd\big)$ 中的极限与余极限逐对象计算可知。
\end{proof}

作为推论，得到：

\begin{othm}{推论}{1.1.5}{n:1.1.5}
函子 $\dR:\PreStk\to\PreStk$ 是函子
\[
\dR|_{\DGSch_{\mathrm{aff}}}:\DGSch_{\mathrm{aff}}\to\PreStk
\]
沿 $\DGSch_{\mathrm{aff}}\hookrightarrow\PreStk$ 的左 Kan 延拓。
\end{othm}

\begin{proof}
对从 $\PreStk$ 到任意 $\infty$-范畴的保余极限函子，这都是对的。
\end{proof}

\ob{1.1.6}{n:1.1.6}此外还有：

\begin{othm}{引理}{1.1.7}{n:1.1.7}
函子 $\dR|_{\DGSch_{\mathrm{aff}}}:\DGSch_{\mathrm{aff}}\to\PreStk$ 同构于函子
\[
\dR|_{{}^{\mathrm{red}}\Sch_{\mathrm{aff}}}:{}^{\mathrm{red}}\Sch_{\mathrm{aff}}\to\PreStk
\]
沿 ${}^{\mathrm{red}}\Sch_{\mathrm{aff}}\hookrightarrow\DGSch_{\mathrm{aff}}$ 的左 Kan 延拓。
\end{othm}

\begin{proof}
对任意目标范畴 $D$ 与任意函子 $\Phi:\DGSch_{\mathrm{aff}}\to D$，映射
\[
\mathrm{LKE}_{{}^{\mathrm{red}}\Sch_{\mathrm{aff}}\hookrightarrow\DGSch_{\mathrm{aff}}}(\Phi|_{{}^{\mathrm{red}}\Sch_{\mathrm{aff}}})\to\Phi
\]
是同构，当且仅当自然变换
\[
\Phi(\clred S)\to\Phi(S),\qquad S\in\DGSch_{\mathrm{aff}}
\]
是同构。取 $D=\PreStk$、$\Phi$ 为函子 $S\mapsto S_{\dR}$ 时，后者按定义成立。
\end{proof}

\ob{1.1.8}{n:1.1.8}设 $\mathcal{C}_1\subset\mathcal{C}_2$ 是 $\PreStk$ 的下列全子范畴中的一对：
\[
{}^{\mathrm{red}}\Sch_{\mathrm{aff}},\ \Sch_{\mathrm{aff}},\ \DGSch_{\mathrm{aff}},\ \DGSch_{\mathrm{qs\text{-}qc}},\ \DGSch,\ \PreStk
\]
（下标「qs-qc」意为「拟分离且拟紧」）。

由引理 1.1.7 与推论 1.1.5 得到：

\begin{othm}{推论}{1.1.9}{n:1.1.9}
由 $\dR|_{\mathcal{C}_2}$ 给出的函子 $\mathcal{C}_2\to\PreStk$，同构于函子 $\dR|_{\mathcal{C}_1}:\mathcal{C}_1\to\PreStk$ 沿 $\mathcal{C}_1\hookrightarrow\mathcal{C}_2$ 的左 Kan 延拓。
\end{othm}

\section{$Y$ 与 $Y_{\dR}$ 的关系}

\ob{1.2.1}{n:1.2.1}函子 $\dR:\PreStk\to\PreStk$ 附带自然变换
\[
p_{\dR}:\mathrm{Id}\to\dR,
\]
即对每个 $Y\in\PreStk$ 有典范映射
\[
p_{\dR,Y}:Y\to Y_{\dR}.
\]

\ob{1.2.2}{n:1.2.2}记 $Y_\bullet/Y_{\dR}$ 为 $p_{\dR,Y}$ 的 Čech 神经，看作 $\PreStk$ 的单纯对象。它以 $Y_{\dR}$ 为增广。

注意每个 $Y_i/Y_{\dR}$ 是 $Y_i$ 沿主对角线的形式完备化。（关于形式完备化的约定，参见 [GR1, Sect.~6.1.1]。）

有典范映射
\begin{equation}
|Y_\bullet/Y_{\dR}|\to Y_{\dR}.\tag{1.2}
\end{equation}

\begin{annote}
直观地说，$Y_i/Y_{\dR}$ 是 $Y^i$ 中「$i$ 个点彼此无穷小接近」所成的形式邻域，而 (1.2) 说的是把这些形式邻域粘起来可以还原出 $Y_{\dR}$。当 $Y$ 经典形式光滑时，这一粘合确实是同构（引理 1.2.4）；一般情形则只是映射。\hfill【来源：[GR1, Sect.~6.1.1]（A）】
\end{annote}

\ob{1.2.3}{n:1.2.3}\textbf{经典形式光滑的预叠。}若对 $S\in\DGSch_{\mathrm{aff}}$，映射
\[
\Maps(S,Y)\to\Maps(\clred S,Y)
\]
在 $\pi_0$ 上诱导满射，则称预叠 $Y$ 经典形式光滑。

下面的结果直接来自定义：

\begin{othm}{引理}{1.2.4}{n:1.2.4}
若 $Y$ 经典形式光滑，则映射
\[
|Y_\bullet/Y_{\dR}|\to Y_{\dR}
\]
是 $\PreStk$ 中的同构。
\end{othm}

\section{局部几乎有限型的情形}

\ob{1.3.1}{n:1.3.1}回忆 $\PreStk$ 含全子范畴 $\PreStk_{\mathrm{laft}}$，由局部几乎有限型的预叠组成，见 [GL:Stacks, Sect.~1.3.9]。按定义，对象 $Y\in\PreStk$ 属于 $\PreStk_{\mathrm{laft}}$，如果：

\begin{itemize}[leftmargin=2.2em]
\item $Y$ 收敛，即对 $S\in\DGSch_{\mathrm{aff}}$，自然映射
\[
\Maps(S,Y)\to\varprojlim_{n\ge 0}\Maps({}^{\le n}S,Y)
\]
是同构，其中 ${}^{\le n}S$ 表示 $S$ 的 $n$-余连接截断；

\item 对每个 $n$，限制 ${}^{\le n}Y:=Y|_{{}^{\le n}\DGSch_{\mathrm{aff}}}$ 属于 ${}^{\le n}\PreStk_{\mathrm{lft}}$，即函子
\[
S\mapsto\Maps(S,Y),\qquad ({}^{\le n}\DGSch_{\mathrm{aff}})^{\mathrm{op}}\to\infty\text{-}\Grpd
\]
与滤过余极限交换（等价地，它是从全子范畴 ${}^{\le n}\DGSch_{\mathrm{aff}}^{\mathrm{ft}}\hookrightarrow{}^{\le n}\DGSch_{\mathrm{aff}}$ 出发的左 Kan 延拓）。
\end{itemize}

\begin{annote}
「收敛」一词在这里是技术性的：它说的是一个预叠的取值由它在各阶截断上的取值决定，正如形式概形由它的无穷小邻域决定。凡本文中出现的 $\PreStk_{\mathrm{laft}}$，都同时要求这两个条件。\hfill【来源：[GL:Stacks, Sect.~1.3.9]（A）】
\end{annote}

\ob{1.3.2}{n:1.3.2}下面的观察在本文中起重要作用。

\begin{othm}{命题}{1.3.3}{n:1.3.3}
设 $Y\in\PreStk_{\mathrm{laft}}$。则：
\begin{enumerate}[label=(\alph*),leftmargin=2.2em,nosep]
\item $Y_{\dR}\in\PreStk_{\mathrm{laft}}$；
\item $Y_{\dR}$ 是经典的，即属于全子范畴 ${}^{\mathrm{cl}}\PreStk\subset\PreStk$。
\end{enumerate}
\end{othm}

\ob{1.3.4}{n:1.3.4}\textbf{(a) 的证明。}

需要验证两条性质：
\begin{enumerate}[label=(\roman*),leftmargin=2.2em,nosep]
\item $Y_{\dR}$ 收敛；
\item 每个截断 ${}^{\le n}(Y_{\dR})$ 局部有限型。
\end{enumerate}
性质 (i) 是同义反复，对任意 $Y\in\PreStk$ 都成立。为证 (ii)，须说明函子 $Y_{\dR}$ 把 ${}^{\le n}\DGSch_{\mathrm{aff}}$ 中的滤过极限送到 $\infty\text{-}\Grpd$ 中的余极限。既然 $Y$ 本身具有这一性质，只需说明函子
\[
S\mapsto\clred S:\DGSch_{\mathrm{aff}}\to\DGSch_{\mathrm{aff}}
\]
保持滤过极限，而这是显然的。

\ob{1.3.5}{n:1.3.5}\textbf{(b) 的证明。}

由推论 1.1.9，须证明余极限
\[
\varinjlim_{S\in(\Sch_{\mathrm{aff}})_{/Y}}S_{\dR}\in\PreStk
\]
是经典的。由 (a)，函子
\[
(\Sch_{\mathrm{aff}}^{\mathrm{ft}})_{/Y}\to(\Sch_{\mathrm{aff}})_{/Y}
\]
共尾；因此
\[
\varinjlim_{S\in(\Sch_{\mathrm{aff}}^{\mathrm{ft}})_{/Y}}S_{\dR}\to\varinjlim_{S\in(\Sch_{\mathrm{aff}})_{/Y}}S_{\dR}
\]
是同构。

于是，由于全子范畴 ${}^{\mathrm{cl}}\PreStk\subset\PreStk$ 对余极限封闭，不妨设 $Y$ 是有限型的经典仿射概形。

更一般地，我们将证明：对 $X\in\DGSch_{\mathrm{aff}}^{\mathrm{aft}}$，预叠 $X_{\dR}$ 是经典的。取闭浸入 $i:X\hookrightarrow Z$，其中 $Z$ 是有限型的光滑经典仿射概形。

记 $Y$ 为 $Z$ 沿 $X$ 的形式完备化 $\widehat{Z}_X$（见 [GR1, Sects.~6.1.1 或 6.5]）。映射 $X\to Y$ 诱导同构 $X_{\dR}\to Y_{\dR}$。因此只需说明 $Y_{\dR}$ 是经典的。

考虑 $Y_\bullet/Y_{\dR}$（见上文 §1.2.2）。注意 $Y$ 形式光滑，因为 $Z$ 形式光滑（见 [GR1, Sect.~8.1]）；特别地，$Y$ 经典形式光滑。由于子范畴 ${}^{\mathrm{cl}}\PreStk\subset\PreStk$ 对余极限封闭，再由引理 1.2.4，只需说明每一项 $Y_i/Y_{\dR}$ 作为预叠是经典的。

注意 $Y_i/Y_{\dR}$ 同构于 $Z^i$ 沿对角嵌入的 $X$ 的形式完备化。于是由 [GR1, Proposition 6.8.2]，$Y_i/Y_{\dR}$ 是经典的。

\begin{annote}
这里的关键一步是把 $X$ 换成它在光滑 $Z$ 中的形式完备化 $Y=\widehat{Z}_X$：$X_{\dR}$ 与 $Y_{\dR}$ 同构，而 $Y$ 形式光滑，于是每一阶无穷小邻域 $Y_i/Y_{\dR}$ 都落在经典概形的范围内。这是 §2.5 中 Kashiwara 引理证明的同一个技巧的原型。\hfill【来源：[GR1, Sect.~6.8]（A）】
\end{annote}

\ob{1.3.6}{n:1.3.6}由命题 1.3.3 得到：

\begin{othm}{推论}{1.3.7}{n:1.3.7}
设 $\mathcal{C}_1\subset\mathcal{C}_2$ 是 $\DGSch_{\mathrm{aff}}$ 的下列全子范畴中的任意一对：
\[
\Sch_{\mathrm{aff}}^{\mathrm{ft}},\quad {}^{<\infty}\DGSch_{\mathrm{aff}}^{\mathrm{aft}},\quad \DGSch_{\mathrm{aff}}^{\mathrm{aft}},\quad \Sch_{\mathrm{aff}},\quad \DGSch_{\mathrm{aff}}.
\]
则对 $Y\in\PreStk_{\mathrm{laft}}$，函子
\[
(\mathcal{C}_1)_{/Y_{\dR}}\to(\mathcal{C}_2)_{/Y_{\dR}}
\]
共尾。
\end{othm}

\begin{proof}
只需对包含
\[
\Sch_{\mathrm{aff}}^{\mathrm{ft}}\hookrightarrow\Sch_{\mathrm{aff}}\hookrightarrow\DGSch_{\mathrm{aff}}
\]
证明论断。对右箭头，断言由命题 1.3.3(b) 得到；对左箭头，由命题 1.3.3(a) 得到。
\end{proof}

\ob{1.3.8}{n:1.3.8}现在考虑下列全子范畴
\begin{equation}
{}^{\mathrm{red}}\Sch_{\mathrm{aff}}^{\mathrm{ft}},\quad \Sch_{\mathrm{aff}}^{\mathrm{ft}},\quad \DGSch_{\mathrm{aff}}^{\mathrm{ft}},\quad \DGSch_{\mathrm{aff}}^{\mathrm{aft}},\quad \DGSch_{\mathrm{laft}},\quad \PreStk_{\mathrm{laft}},\tag{1.3}
\end{equation}
它们是 §1.1.8 中出现的那些范畴的子范畴。

\begin{othm}{推论}{1.3.9}{n:1.3.9}
函子 $\dR$ 到 $\PreStk_{\mathrm{laft}}$ 上的限制，同构于它到 $\mathcal{C}$ 上的左 Kan 延拓，其中 $\mathcal{C}$ 是 (1.3) 中的任一子范畴。
\end{othm}

\begin{proof}
只需对 $\mathcal{C}={}^{\mathrm{red}}\Sch_{\mathrm{aff}}^{\mathrm{ft}}$ 证明。由推论 1.1.9，只需说明对 $Y\in\PreStk_{\mathrm{laft}}$，函子
\[
({}^{\mathrm{red}}\Sch_{\mathrm{aff}}^{\mathrm{ft}})_{/Y}\to(\Sch_{\mathrm{aff}})_{/Y}
\]
共尾。

由命题 1.3.3(a)，函子
\[
(\Sch_{\mathrm{aff}}^{\mathrm{ft}})_{/Y}\to(\Sch_{\mathrm{aff}})_{/Y}
\]
共尾。再由包含 ${}^{\mathrm{red}}\Sch_{\mathrm{aff}}^{\mathrm{ft}}\hookrightarrow\Sch_{\mathrm{aff}}^{\mathrm{ft}}$ 有右伴随这一事实，即得断言。
\end{proof}
